\documentclass[11pt]{article}
\usepackage{fontspec}
\usepackage{polyglossia}
\setmainlanguage{russian}
\setotherlanguage{english}
\setmainfont{DejaVu Serif}
\setmonofont{DejaVu Sans Mono}
\usepackage{amsmath}
\usepackage{booktabs}
\usepackage{graphicx}
\usepackage{hyperref}
\usepackage{geometry}
\geometry{margin=1in}
\hypersetup{colorlinks=true, linkcolor=black, citecolor=black, urlcolor=blue}

\title{Извлечение кортежей мнений из русскоязычных новостей:\\
дообученная русскоязычная LLM против многоязычных решений}
\author{Грузин Н.А. Б13-303}
\date{}

\begin{document}
\maketitle

\begin{abstract}
Извлечение мнений из текста редко сводится к одному ярлыку тональности: чтобы сигнал был полезен,
приходится знать, кто высказывается, о ком и какими именно словами. Мы решаем эту структурную постановку
в формулировке соревнования RuOpinionNE-2024, где из новостного предложения требуется восстановить все
кортежи (источник, объект, выражение, тональность). Сущности здесь почти всегда
организации и персоны, поэтому извлечённый кортеж работает как entity-level сигнал тональности, удобный, например,
при анализе новостного фона по отдельным компаниям. Опираясь на русскоязычную модель T-pro-it-2.0 (Qwen3-32B) и дообучив
её методом QLoRA с генерацией строгого JSON и последующей агрегацией нескольких сэмплов, мы получаем
F1 = 0.460 на единственном публично размеченном наборе. На том же наборе это превосходит и лучшую систему
фазы валидации (0.34), и базовое решение организаторов (0.17); полученное число выше и лучшего результата
фазы теста (0.41, дообученная LLaMA-3.3-70B), хотя валидация и тест это разные сплиты, поэтому последнее
сравнение прямым не является. Победитель шёл схожим путём, дообучая модель и агрегируя несколько ответов,
так что наш интерес не в самой схеме, а в том, что русскоязычной модели вдвое меньшего размера хватает для более
высокого результата на публичной валидации.
Код и все конфигурации эксперимента открыты.
\end{abstract}

\section{Введение}
Новостной поток остаётся одним из самых ранних источников информации о компаниях и рынках, и для прикладного
анализа важна не усреднённая окраска заметки, а отношение к конкретному эмитенту: кто его высказал, к кому оно
обращено и чем подкреплено. Обычная классификация тональности предложения такого ответа не даёт, теряя ровно ту структуру,
ради которой текст и читают.

Восстанавливая эту структуру, задача structured sentiment analysis превращает предложение в набор кортежей
мнений. Для русского языка её задаёт соревнование RuOpinionNE-2024~\cite{ruopinionne2024}, продолжающее
RuSentNE-2023~\cite{rusentne2023}. Поскольку размечены в этих новостях преимущественно организации и персоны,
аккуратно извлечённый кортеж оказывается тем самым entity-level сигналом, что используется в финансовом NLP,
но дополненным указанием источника оценки и её словесного обоснования.

Задача трудная: лучшее решение соревнования остановилось на 0.41 F1, базовое набрало 0.24 на тесте и 0.17 на
валидации, так что запас для улучшения заведомо велик. Победитель добивался своего числа дообучением
LLaMA-3.3-70B и агрегацией нескольких генераций, поэтому сама по себе схема, то есть дообучение плюс голосование,
не нова. Наш вопрос ставится иначе: хватит ли русскоязычной T-pro-it-2.0 на базе Qwen3-32B, вдвое меньшей по
числу параметров, чтобы под той же схемой получить более высокий результат. Как выясняется, хватает. Наш вклад:
\begin{itemize}
\item постановка задачи и разбор метрики, доведённый до краевых случаев разметки, на которых обычно и теряются баллы;
\item дообученное QLoRA-решение на T-pro-it-2.0, превосходящее лучший опубликованный результат фазы валидации при вдвое меньшем размере модели;
\item сравнение с конкурентами по единой метрике: та же модель без дообучения, вклад self-consistency и опубликованные участники;
\item воспроизводимый пайплайн с открытым кодом, считающий качество официальным скриптом.
\end{itemize}

\section{Обзор работ}
\paragraph{Structured sentiment analysis.}
Извлечение кортежей мнений, восходящее к SemEval-2022 Task~10~\cite{barnes2022semeval}, изначально ставилось
как разметка графа, в котором узлами служат спаны источника, объекта и выражения, а рёбрами полярные связи между
ними. Ранние системы, опираясь на разбор зависимостей и на энкодеры с CRF, строили этот граф напрямую. С приходом
инструктивных LLM задачу всё чаще переформулируют как генерацию: модель печатает кортежи текстом, после чего спаны
выравнивают по исходному предложению.

\paragraph{Целевая тональность для русского.}
RuSentNE-2023~\cite{rusentne2023} решал более узкую задачу, спрашивая лишь тональность предложения к заранее
указанной сущности. В развивающей эту линию работе~\cite{thor2024} дообученная Flan-T5-xl, рассуждавшая
по схеме three-hop chain-of-thought (THoR), обошла лучший энкодерный результат соревнования, и здесь видна общая
закономерность: рассуждение перед ответом помогает целевой тональности. RuOpinionNE-2024~\cite{ruopinionne2024},
добавив извлечение источника и обоснования, поднимает планку до полного кортежа.

\paragraph{Результаты RuOpinionNE-2024.}
Обзор соревнования~\cite{ruopinionne2024} приводит две таблицы. На фазе валидации лучшим оказался результат
0.34 F1 (Zholtikov\_Michail), тогда как базовое решение, то есть Qwen2.5-32B в режиме 10-shot, дало лишь 0.17.
На финальном тесте победила дообученная в 4-бит LLaMA-3.3-70B с агрегацией пяти ответов (VatolinAlexey, 0.41),
за ней расположились 0.35, 0.33 и 0.28, а базовое решение набрало 0.24. Победитель, как видно, сделал ставку
на крупную англоцентричную модель и на self-consistency; мы сохраняем идею агрегации, но меняем основу на
русскоязычную модель меньшего размера.

\paragraph{LLM для русского.}
Линейка T-lite/T-pro~\cite{tpro2025}, адаптирующая Qwen2.5/Qwen3 под русский переобучением токенизатора и
продолженным предобучением, к 2025 году заметно подтянула качество на русскоязычных бенчмарках; построенная
на Qwen3-32B модель T-pro-it-2.0 как раз из этого семейства. Чтобы уместить её 32 миллиарда параметров в
4-бит на одной A100, мы дообучаем её методом QLoRA~\cite{qlora2023}, надстроенным над LoRA~\cite{lora2021}.

\paragraph{Финансовый NLP.}
Тональность в финансах традиционно снимали моделями уровня FinBERT~\cite{finbert2019}, обученными на наборах
вроде Financial PhraseBank~\cite{phrasebank2014}. Давая оценку всему сообщению целиком, такие модели не отвечают
на вопрос, к кому именно она относится. Извлечение кортежей мнений этот пробел закрывает, сразу указывая объект
оценки и её обоснование, то есть именно ту информацию, которой пословная классификация не даёт.

\paragraph{Self-consistency.}
Голосование по нескольким сэмплам генерации~\cite{selfconsistency2022}, предложенное изначально для арифметических
рассуждений, мы переносим на извлечение: кортеж принимается, если он встретился не менее чем в $k$ из $N$ независимых
генераций. Снижая разброс, такая агрегация заодно отсекает редкие галлюцинации.

\section{Данные и постановка задачи}
\paragraph{Формат.}
На входе подаётся предложение из новостного текста, на выходе ожидается множество кортежей мнений, каждый из
которых состоит из четырёх частей. Источником мнения выступает либо именованная сущность, либо специальная метка
\texttt{AUTHOR} (когда оценку высказывает автор текста), либо \texttt{NULL} (когда источник не выражен); объектом
служит сущность; выражение это один или несколько фрагментов, обосновывающих оценку; тональность принимает значения
\texttt{POS} или \texttt{NEG}. Источник, объект и выражение задаются точными подстроками предложения с символьными
смещениями.

\paragraph{Объём и распределение.}
Статистика наборов сведена в таблицу~\ref{tab:data}, и в ней обращают на себя внимание две особенности. Во-первых,
около 41\% предложений вовсе не содержат мнений, а значит, модель обязана уметь возвращать пустой ответ, иначе
страдает точность. Во-вторых, число кортежей на предложение распределено длиннохвосто: чаще всего их ноль, один
или два, но попадаются предложения и с десятком, и даже с двадцатью тремя кортежами. По типу источника преобладает
именованная сущность, заметную долю составляют мнения без выраженного источника (\texttt{NULL}), и примерно
четырнадцать процентов исходят от автора. Отдельную трудность создают разрывные выражения, обоснование в которых
собрано из несоседних фрагментов (например, "снял \dots\ ограничения на реализацию"). Утечки между обучением и
валидацией нет: ни одно из 1316 валидационных предложений в обучающей части не встречается, что мы проверили явно.

\begin{table}[h]
\centering
\caption{Статистика RuOpinionNE-2024. Размер теста дан по выложенному файлу \texttt{test.jsonl}
(803 предложения); в обзоре~\cite{ruopinionne2024} приводится 804.}
\label{tab:data}
\begin{tabular}{lccc}
\toprule
 & Train & Validation & Test \\
\midrule
Предложений & 2556 & 1316 & 803 \\
\quad из них без мнений & 1062 & 547 & n/a \\
Кортежей & 2904 & 1612 & скрыто \\
\quad POS\,/\,NEG & 1174\,/\,1730 & 753\,/\,859 & n/a \\
Источник entity\,/\,null\,/\,author & 1623\,/\,897\,/\,384 & 980\,/\,436\,/\,196 & n/a \\
Разрывных выражений & 19 & 6 & n/a \\
\bottomrule
\end{tabular}
\end{table}

\paragraph{Метрика.}
Качество измеряется взвешенным tuple-F1 (официальный скрипт соревнования). Предсказанный кортеж засчитывается
совпавшим с эталонным тогда, когда источник, объект и выражение пересекаются по токенам хотя бы на один токен,
а тональность совпадает точно. Поскольку вклад совпадения взвешивается средней долей пересечения по трём спанам,
частичные попадания учитываются мягко, не обнуляясь. Точность и полнота считаются по всему набору и затем сводятся
в F1, причём метка \texttt{NULL} в роли источника трактуется отдельно от \texttt{AUTHOR}.

\paragraph{Протокол оценки.}
Соревнование на Codalab закрыто, а разметка теста не опубликована, поэтому единственным публично размеченным
набором остаётся валидация. Обучаясь только на обучающей части и оценивая результат на валидации официальным
скриптом, мы получаем число, прямо сопоставимое с таблицей фазы валидации (лучшая система 0.34, базовое решение
0.17). Числа фазы теста мы приводим для полноты, оговаривая, что шкалы фаз заметно различаются.

\paragraph{Прикладной сценарий.}
Сама задача шире финансов, но один из естественных сценариев применения это финансовые новости, где
сущности это в основном компании, банки, регуляторы и должностные лица. Извлечённый кортеж
(источник, объект, выражение, тональность) даёт здесь entity-level сигнал,
отвечающий, к какой компании относится оценка, кто её высказал и чем она обоснована, и, в отличие от пословной
тональности, допускает агрегацию по конкретной компании во времени.

\section{Подход}
\paragraph{Постановка как генерация.}
Формулируя извлечение как инструктивную генерацию, мы описываем в системном промпте определение кортежа и требуем
на выходе строгий JSON-массив объектов с полями \texttt{holder}, \texttt{target}, \texttt{expression} (список
фрагментов) и \texttt{polarity}. Печатая подстроки исходного предложения, модель не выдаёт смещения напрямую,
поэтому символьные позиции мы восстанавливаем поиском подстроки и приводим к официальному формату. Корректность
этого преобразования подтверждена round-trip тестом: эталон, прогнанный через сериализацию, парсинг и обратное
восстановление смещений, даёт F1 = 0.999, а значит, формат и парсер информацию практически не теряют.

\paragraph{Санитизация.}
Официальный скрипт прерывается с ошибкой, если в одном предложении два кортежа делят источник, объект и
тональность при пересекающихся выражениях. Чтобы предсказание оставалось валидным при любом выходе модели, мы
пропускаем его через дедупликацию, повторяющую логику официального конвертера, но отбрасывающую конфликтный
кортеж вместо того, чтобы упасть.

\paragraph{Основная модель.}
Дообучая T-pro-it-2.0 (Qwen3-32B) методом QLoRA, мы держим базовые веса в 4-бит (nf4, double quantization,
вычисления в bfloat16) и обучаем LoRA-адаптеры на всех линейных проекциях (ранг 32, $\alpha = 64$, dropout 0.05).
Лосс считается только по токенам ответа, тогда как промпт маскируется. Оптимизация ведётся paged AdamW 8-bit при
learning rate $10^{-4}$ с косинусным спадом и прогревом в 5\%; обучение длится три эпохи при эффективном батче 16
и ограничении длины в 2048 токенов, с включённым gradient checkpointing. Режим рассуждения отключён намеренно:
печатая JSON сразу, модель экономит токены и время инференса.

\paragraph{Self-consistency.}
На инференсе мы порождаем $N$ независимых сэмплов при ненулевой температуре и сводим их голосованием, принимая
кортеж, если он встретился не менее чем в $k$ из $N$ ответов; порог $k$ подбирается по валидации. При $N=1$ это
вырождается в обычное жадное декодирование, тогда как голосование, снижая разброс на небольшом наборе, убирает
редкие нестабильные кортежи.

\paragraph{Конкурентные подходы.}
Сравнение мы ведём по трём осям. Первая это та же T-pro-it-2.0 без дообучения, в режимах zero-shot и few-shot,
позволяющая отделить вклад дообучения от качества самой модели. Вторая это вклад self-consistency, измеряемый
сопоставлением жадного декодирования с агрегацией сэмплов. Третья это опубликованные результаты участников,
включая решение-победитель на LLaMA-3.3-70B. Наши собственные числа посчитаны одним официальным скриптом на одном
наборе, поэтому сопоставимы между собой без оговорок; числа участников фазы теста, как отмечено выше, относятся к
другому закрытому сплиту и приводятся лишь ориентировочно.

\section{Эксперименты и результаты}
\paragraph{Окружение.}
Все запуски выполнены на одной NVIDIA A100 80GB; инференс вёлся через vLLM с загрузкой базовой модели в 4-бит,
а качество измерялось официальным скриптом RuOpinionNE-2024 на валидационном наборе.

\begin{table}[h]
\centering
\caption{Результаты на валидационном наборе RuOpinionNE-2024 (взвешенный tuple-F1).
Числа предыдущих работ взяты из обзора соревнования~\cite{ruopinionne2024}.}
\label{tab:main}
\begin{tabular}{lcc}
\toprule
Система & Размер & F1 (валидация) \\
\midrule
Базовое решение (Qwen2.5-32B, 10-shot)~\cite{ruopinionne2024} & 32B & 0.17 \\
Лучшая система фазы валидации~\cite{ruopinionne2024} & n/a & 0.34 \\
\midrule
T-pro-it-2.0, zero-shot (наша) & 32B & 0.191 \\
T-pro-it-2.0, few-shot 3 (наша) & 32B & 0.199 \\
T-pro-it-2.0 + QLoRA, жадное декодирование (наша) & 32B & 0.438 \\
\textbf{T-pro-it-2.0 + QLoRA + self-consistency (наша)} & \textbf{32B} & \textbf{0.460} \\
\bottomrule
\end{tabular}
\end{table}

\noindent
Наш результат 0.460 выше опубликованного лучшего результата фазы валидации (0.34) и базового решения (0.17)
на том же наборе. Для полноты приведём и числа фазы теста~\cite{ruopinionne2024}: победитель (LLaMA-3.3-70B с
агрегацией) набрал 0.41, далее идут 0.35 и 0.33, базовое решение 0.24. Эти числа относятся к закрытому тесту на
другом сплите, поэтому прямого сравнения с нашим валидационным результатом они не дают и приводятся лишь
ориентировочно.

\paragraph{Вклад дообучения.}
Без дообучения T-pro набирает всего 0.191 в zero-shot и 0.199 в few-shot: правильную JSON-структуру модель печатает,
однако границы спанов угадывает плохо и часть мнений упускает. Примеры в контексте почти не выручают, и это
объяснимо, ведь задача упирается не в формат, а в точное выделение источника, объекта и выражения. Дообучение же
поднимает результат до 0.438, из чего следует, что основной прирост даёт адаптация под конкретную разметку, а не
размер модели и не few-shot.

\paragraph{Выбор эпохи.}
Сохраняя адаптер после каждой эпохи, мы оценили все три (таблица~\ref{tab:epoch}). Прирост между второй и третьей
эпохами невелик, а переобучения на валидации не наблюдается, поэтому в качестве рабочей мы берём третью.

\begin{table}[h]
\centering
\caption{F1 по эпохам (жадное декодирование).}
\label{tab:epoch}
\begin{tabular}{cccc}
\toprule
Эпоха & 1 & 2 & 3 \\
\midrule
F1 & 0.395 & 0.435 & 0.438 \\
\bottomrule
\end{tabular}
\end{table}

\paragraph{Self-consistency.}
Порождая пять сэмплов при температуре 0.6, мы принимаем кортеж по порогу $k$ из пяти, и этот порог регулирует
баланс полноты и точности. При $k=1$ берётся объединение всех сэмплов (F1 = 0.435), при $k=5$ остаётся их
пересечение (0.299), а оптимум достигается на $k=2$, давая 0.460 и прибавляя 0.022 к жадному декодированию.
Отбрасывая нестабильные кортежи, голосование при чрезмерно высоком пороге начинает терять и верные, чем и
объясняется спад на $k\ge 3$. Голосовать можно двумя способами: по точному совпадению нормализованных строк и
с кластеризацией кандидатов по тому же критерию пересечения токенов, что заложен в метрику. Оптимум у обоих близок
(0.460 и 0.461 при $k=2$), однако кластеризация заметно устойчивее к порогу: при $k=5$ она удерживает 0.39 против
0.30 у строкового варианта, поскольку перестаёт штрафовать кортежи лишь за расхождение в границах спанов.

\paragraph{Анализ ошибок.}
На лучшем прогоне модель нашла 737 кортежей при 871 пропуске и 750 ложных срабатываниях, так что основная потеря
приходится на полноту: распознаётся примерно половина эталонных кортежей. Тональность при этом почти не путается
(всего шесть случаев, когда спаны совпали, а знак нет), а значит, найдя кортеж, модель определяет его полярность
надёжно. Тяжелее всего даются мнения без явного источника (\texttt{NULL}), где доля пропусков максимальна; разрывные
же выражения, будучи редкими, на итог почти не влияют. Главным резервом роста остаётся полнота на неявно выраженных
мнениях.

\section{Заключение}
Решая задачу извлечения кортежей мнений из русских новостей, мы получили на единственном публично размеченном
наборе F1 = 0.460, что выше опубликованного лучшего результата фазы валидации (0.34) и базового решения (0.17).
Результат победителя (0.41) получен на закрытом тесте, то есть на другом сплите, поэтому прямым сравнением с нашим
числом он не является. Определяющим оказался выбор основы: при той же схеме из дообучения и self-consistency
русскоязычной T-pro-it-2.0 на 32 миллиарда параметров хватило, чтобы превзойти валидационный результат соревнования,
оставаясь вдвое меньше модели-победителя. В приложении к финансам извлечённые кортежи дают
entity-level сигнал тональности к компаниям, пригодный для агрегации новостного фона по отдельному эмитенту, а
основным направлением дальнейшей работы видится повышение полноты на мнениях без явного источника.

\begin{thebibliography}{99}
\bibitem{ruopinionne2024} Lukashevich N., Rusnachenko N. et al. RuOpinionNE-2024: Extraction of Opinion Tuples from Russian News Texts. arXiv:2504.06947, 2025.
\bibitem{rusentne2023} Golubev A., Rusnachenko N., Loukachevitch N. RuSentNE-2023: Named Entity Oriented Sentiment Analysis. arXiv:2305.17679, 2023.
\bibitem{thor2024} Rusnachenko N., Loukachevitch N. Large Language Models in Targeted Sentiment Analysis for Russian. arXiv:2404.12342, 2024.
\bibitem{barnes2022semeval} Barnes J. et al. SemEval-2022 Task 10: Structured Sentiment Analysis. SemEval, 2022.
\bibitem{tpro2025} T-Tech. T-pro-it-2.0 (Qwen3-32B, russian-adapted). Hugging Face, 2025.
\bibitem{qlora2023} Dettmers T. et al. QLoRA: Efficient Finetuning of Quantized LLMs. arXiv:2305.14314, 2023.
\bibitem{lora2021} Hu E. et al. LoRA: Low-Rank Adaptation of Large Language Models. arXiv:2106.09685, 2021.
\bibitem{selfconsistency2022} Wang X. et al. Self-Consistency Improves Chain of Thought Reasoning. arXiv:2203.11171, 2022.
\bibitem{finbert2019} Araci D. FinBERT: Financial Sentiment Analysis with Pre-trained Language Models. arXiv:1908.10063, 2019.
\bibitem{phrasebank2014} Malo P. et al. Good Debt or Bad Debt: Detecting Semantic Orientations in Economic Texts. JASIST, 2014.
\end{thebibliography}

\end{document}
